\part{Coherent states}\label{part:coherent}
\section{Coherent overlap and a uniform trace moment}
\label{rec:sec-coherent}

We next ask how the overlap changes when missing positions are retained as an ordered list. Write $n=2^d$ and $B_n=K_d$ for $d\ge1$; for $d=0$
put $B_1=\id$, the empty sweep. The method proves a dimension-weighted trace bound with right side one.  It combines a partial-mapping density
estimate with a coherent-state overlap on a grid with holes.  The
coherent overlap uses a moment-matching change of basis; we include
that step because it explains why the grid factors have mean at most
one even after cells are removed.

\begin{theorem}\label{rec:coherent-main}
There is an absolute $p_*\ge1$ such that, for every dyadic $n$ and
every irreducible type $\lambda\vdash n$,
\begin{equation}\label{rec:coherent-unit-moment}
 D_\lambda\Tr\!\left(((B_n^*B_n)|_{V_\lambda})^{p_*}\right)\le1.
\end{equation}
Consequently $\|B_n|_{V_\lambda}\|\le D_\lambda^{-1/(2p_*)}$. The moment in the display is the Schatten power $2p_*$, because $\operatorname{Tr}(B_n^*B_n)^{p_*}=\|B_n\|_{2p_*}^{2p_*}$.
\end{theorem}

\subsection{Two sparse estimates}

For a nontrivial $\lambda=(n-k,\beta)$ with $k=n-\lambda_1\ge1$, the two inputs are
\begin{align}
 \|B_n|_{V_\lambda}\|^2
 &\le C^k(1+d)^{Ck}(k/n)^{k/2},
       \label{rec:coherent-contact-bound}\\
 \|B_n|_{V_\lambda}\|^2
 &\le e^{Ck}D_\beta^{-a}
       \label{rec:coherent-tail-bound}
\end{align}
for absolute $a>0$.
For the first estimate we use the square of the fully cancelled kernel. Work with counting-measure matrices on ordered injective $k$-tuples. For $A\subseteq[k]$, let $K_A^{\rm prob}$ be the sweep transition kernel on the coordinates in $A$, and set
\[
 L_A(x,y)=n^{-(k-|A|)}K_A^{\rm prob}(x_A,y_A),\qquad
 D(x,y)=\sum_{A\subseteq[k]}(-1)^{k-|A|}L_A(x,y).
\]
Both arguments of $L_A$ are restricted to injective full tuples. Its row and column sums are at most one: the coordinates in $A$ follow a genuine sweep and the others are sampled independently uniformly, after which noninjective outcomes are discarded. Consider the subspace of functions whose sum over any one tuple
coordinate is zero when the other coordinates are fixed.  All
proper-coordinate marginals vanish on this subspace, so $D$ acts
exactly as the full sweep kernel $K=L_{[k]}$.  The type
$\lambda=(n-k,\beta)$ belongs to it by the branching argument in
Section~\ref{ten:weighted:part:21}.

The product-of-pair-factors formula for the sweep shows that $D(x,y)=0$ if one of the $k$ unique start-to-end paths has no potential switch contact with another path. Write $\mathcal E(x,y)$ for the event that this contact graph has no isolated vertex. By discarding this condition on the first leg only,
\[
 |(D^*D)(x,y)|
 \le\sum_{A,B\subseteq[k]}\ \sum_{z\text{ injective}}
       L_A(z,x)L_B(z,y)\mathbf1_{\mathcal E(z,y)}.
\]
For fixed $x,A$, the first-leg mass in $z$ is bounded by the law obtained from a reverse sweep of the labels in $A$ and independent uniform positions for the others. Under this unrestricted mixed law, the occupied set $S$ satisfies
\[
 \Pr(F\subseteq S)\le(k/n)^{|F|}
\]
for every fixed site set $F$. For the sweep subset this follows from the negative-inclusion argument in \eqref{ten:weighted:eq:2}. Adding the independent uniform positions preserves the bound: partition the tested sites between the two sources and sum the resulting product bounds. Restriction to injective $z$ only decreases these probabilities; it is retained in the kernel sum, so the relevant $S$ always has exactly $k$ sites.

Independently presample a complete common routing for the second leg and one private uniform path from each site. Any assignment of the labels in $B$ to common paths and the other labels to private paths is covered by the $2^k$ choices of path type on a $k$-site set. Thus no factor $k!$ is needed. The expected number of such typed unordered sets with no isolated contact is at most
\begin{equation}\label{rec:coherent-forest-count}
 \sum_{1\le j\le k/2}C^k\binom nj(Cd)^{k-j}.
\end{equation}
Indeed choose roots, ordered forest shapes and path types. Every nontrivial component has a rooted spanning tree. A new common-path child has at most two possible starts per switch; summing over private-path starts gives at most two expected choices per switch. Expose independent paths only at unused sites and sum successively over the tree extensions. This bounds each extension by $4d$ and proves the display, with the number of forest shapes absorbed into $C^k$.

Multiply the expected count by $(k/n)^k$, using the first-leg inclusion bound independently of the second-leg fields, and sum the at most $4^k$ choices of $A,B$. The elementary bound $\binom nj\le(Cn/k)^{k/2}$ for $1\le j\le k/2$ shows that every absolute row sum of $D^*D$ is at most
\[
 R_k=C^k(1+d)^{Ck}(k/n)^{k/2}.
\]
For $k=1$ the contact event is empty and the same assertion is immediate. Since $D^*D$ is self-adjoint, the column sums obey the same bound, and the Schur test gives $\|D^*D\|\le R_k$. On that zero-marginal subspace, $D=K$, so $\|K\|^2\le R_k$. This proves \eqref{rec:coherent-contact-bound} without another square-root loss.

The second estimate is Proposition~6.3 of \citet{OpenAI2026Routing}, applied with $T_N=B_n$ and $N=n$. Its input, Theorem~6.1 of that companion, is the stronger density statement
\begin{equation}\label{rec:coherent-density}
 \sup_x n^{-k}\sum_y
       \{n^k\Pr_{B_n^*B_n}(x\mapsto y)\}^{1+\delta}
       \le e^{Ck}
\end{equation}
for a fixed absolute $0<\delta\le1$, all $0\le k\le n$, and every ordered injection $x$. Repeated-site output tuples contribute zero. The companion gives the displayed power normalization as well as the normalization by $(n)_k$, in both orientations and for both rows and columns. Its Proposition~6.3 transfers this bound through the right $S_k$ label action to obtain exactly \eqref{rec:coherent-tail-bound}, with $a=\delta/(1+\delta)$. Section~6.6 there also proves the density estimate by the distinct cycle-window certificate and entropy argument. Only these partial-deck statements are used here; the moment induction below remains local.

For sufficiently large $d$, put $s=\log(n/k)$ and suppose
$s\ge d^{3/4}$.  The first bound then gives
$\log\|B_n|_{V_\lambda}\|^2\le-ks/4$.
Taking the geometric mean of this bound and
\eqref{rec:coherent-tail-bound}, and increasing the fixed lower
threshold on $d$, gives
\[
 \log\|B_n|_{V_\lambda}\|^2
 \le-ks/8+Ck/2-(a/2)\log D_\beta
 \le-ks/16-(a/2)\log D_\beta.
\]
Since $\log D_\lambda\le k(1+s)+\log D_\beta$, there is an
absolute $c>0$ such that $\|B_n|_{V_\lambda}\|^2\le D_\lambda^{-c}$
in this range.  In particular,
\[
 D_\lambda\Tr(B_n^*B_n)^p
 \le D_\lambda^2\|B_n|_{V_\lambda}\|^{2p}\le1
 \qquad(p\ge2/c).
\]
This supplies the trace moment directly for the sparse diagrams.

\subsection{A coherent overlap on a punctured grid}

We first estimate the overlap before restoring the holes.  The
projections below select a single carrier direction in each line,
while retaining every multiplicity copy.

\begin{proposition}[Coherent overlap on surviving cells]
\label{prop:coherent-hook}
Let $M,R,Q\ge1$ be integers.  Consider an $M$ by $R$ grid,
$n=MR$, with $l$ holes and $m=n-l$ surviving cells.  Let $K,J$ be the row and column permutation groups
on surviving cells.  Suppose $\mu\vdash m$ lies in a $(Q,Q)$ hook.
For Fourier projections $E=\bigotimes_i E_{\sigma_i,e_i}$ and $F=\bigotimes_j E_{\tau_j,f_j}$ selecting one unit carrier vector on each row or column, write $\sigma=(\sigma_i)$ and $\tau=(\tau_j)$. Then
\begin{equation}\label{rec:coherent-hook-grid}
 D_\mu\Tr_{V_\mu}(EF)
 \le D_\sigma D_\tau
       \exp\{l+C(1+M+R)Q^2\log(2n)\}.
\end{equation}
Here $D_\sigma,D_\tau$ denote products over lines, and $C$ is
absolute.  For $m=0$ the
surviving representation is the one-dimensional trivial representation.
\end{proposition}

\begin{proof}
The case $m=0$ is immediate.  Assume $m>0$ so the count
proportions below are defined.  Use the signed alphabet with $Q$ even and $Q$ odd symbols.
For the compact group $U(Q)\times U(Q)$ its tensor representation
decomposes into carriers $W_{\alpha\beta}$ paired with
$\operatorname{Ind}(V_\alpha\otimes V_{\beta^{\mathsf T}})$.
The carrier dimensions and the multiplicities of compatible types
are at most $e^{CQ^2\log(2n)}$.  For a probability vector $q$, write
$H(q)=-\sum_a q_a\log q_a$, with $0\log0=0$; for a diagonal
state $\Lambda$, $H(\Lambda)$ means the entropy of its diagonal.
If $V_\gamma$ occurs in the multiplicity space paired with
$(\alpha,\beta)$, put $c=(\alpha,\beta)$ and $u=|\gamma|$.
The hook formula gives, with $uH(c/u)=0$ when $u=0$,
\begin{equation}\label{rec:coherent-count-dimension}
 D_\gamma\ge e^{uH(c/u)-CQ^2\log(2n)}.
\end{equation}
Indeed the first $Q$ row lengths differ from $\alpha_i$ by at most
$Q$, and the first $Q$ lower column lengths differ from $\beta_i$
by at most $Q$.  Their hook product differs from the corresponding
row and column factorials by at most the displayed error.
Every $(Q,Q)$-hook diagram occurs by taking its first $Q$ rows and
the transpose of the remainder.

Choose a compact-group sector whose multiplicity space contains
$V_\mu$.  Let $c_*$ be its highest-weight count and put
$\Lambda=\operatorname{diag}(c_*/m)$.
Then $D_\mu\le e^{mH(\Lambda)}$.  Let $P$ project onto that highest
line in every copy of its compact type.  On the signed slot space
$\mathcal L_m=(\mathbb C^Q\oplus\mathbb C^Q)^{\otimes m}$,
$P$ commutes with the position-permutation action, and its range
contains $V_\mu$.  Since the contribution of each irreducible to
the trace of a product of orthogonal projections is nonnegative,
\[
 \Tr_{V_\mu}(EF)\le\Tr_{\mathcal L_m}(E P F P).
\]
We will resolve $E,F$ into positive rank-one terms.  The last trace
then becomes a positive weighted sum of $|\langle y,Px\rangle|^2$.
We explain which factor is averaged in the coherent-orbit resolution \cite[Section~2, equations~(7)--(11)]{Perelomov1972}. On a line with $u$ surviving slots the signed tensor space has the orthogonal decomposition
\[
 \mathcal L_u=\bigoplus_\kappa W_\kappa\otimes M_\kappa,
 \qquad
 M_\kappa=\bigoplus_\gamma V_\gamma\otimes\mathcal N_{\kappa,\gamma},
\]
where $W_\kappa$ is a compact-group carrier of highest count $c_\kappa$, and $M_\kappa$ carries the induced symmetric-group representation. A matrix action selecting a unit vector $e$ in type $\gamma$ acts on this summand as
\[
 I_{W_\kappa}\otimes|e\rangle\langle e|
                  \otimes I_{\mathcal N_{\kappa,\gamma}}.
\]
With $dU$ the normalized Haar probability measure on
$U(Q)\times U(Q)$ and $h_\kappa$ a unit highest vector, replace
the first identity by
\[
 I_{W_\kappa}=\dim W_\kappa
       \int|Uh_\kappa\rangle\langle Uh_\kappa|\,dU,
\]
and resolve the last identity in an orthonormal basis. This gives positive rank-one terms whose vectors, in the slot realization, are unitary translates of highest-count vectors. The arbitrary selected Specht vector $e$ remains unchanged; all multiplicity copies are present. The dimensions of the compact carriers, the number of compatible sectors and the multiplicities are at most $e^{CQ^2\log(2n)}$, so the total weight costs this factor per line. There is no factor $D_\gamma$, because the selected state $|e\rangle\langle e|$ has trace one. Tensoring these resolutions across the lines gives the row and column coherent product vectors used next. Such a vector can be entangled within a line; ``product'' here refers to different lines.

It remains to bound a coefficient $\langle y,Px\rangle$, where
$x,y$ are unit coherent product vectors on rows and columns.  Let
$\theta_i,\phi_j$ be their normalized line count vectors, with line
sizes $r_i,s_j$.  For an empty line choose any probability vector;
its entropy contribution has coefficient zero.  We claim
\begin{equation}\label{rec:coherent-coefficient}
 |\langle y,Px\rangle|^2
 \le\exp\left\{CQ^2\log(2n)-mH(\Lambda)
          +\sum_ir_iH(\theta_i)+\sum_js_jH(\phi_j)+l\right\}.
\end{equation}

To control the product over surviving cells, we will arrange that
the row and column count densities have the same average
$\Lambda$.  The required change of basis must also preserve the
coefficient up to factors at most one.  A minimum-norm argument
provides both properties.  This viewpoint has its classical
antecedent in Kempf--Ness \cite[Theorem~0.1]{KempfNess1979};
for a character-normalized highest-weight formulation, compare
\citet[Sections~2 and~4]{FranksWalter2022}. We give the local argument because the occupied-cell and coefficient-square bounds require its exact normalization.  Let $g$ be
block lower triangular on the even and odd alphabets, acting on
the slot space by $g^{\otimes m}$.  With the highest-weight
convention used here, projection onto the highest line satisfies
$Pg^{\otimes m}=\chi(g)P$, where
$\chi(g)=\operatorname{diag}(g)^{c_*}$.  We suppress the tensor-power
notation in the orbit argument.
We justify the limiting normalization explicitly. If $Px=0$ or $Py=0$, the coefficient is zero and no scaling is needed. Assume both are nonzero. Consider the character-corrected orbit
\[
 \mathcal O_x=\{\chi(g)^{-1}g x:g\text{ is invertible and block lower triangular}\}.
\]
Its norm infimum $c_x$ satisfies $0<\|Px\|\le c_x\le1$, because every vector in the orbit has projection $Px$ and the identity is allowed. Choose a norm-minimizing sequence. It is bounded, so a subsequence converges to a vector $z_x$ with norm $c_x$ in the orbit closure. That closure is invariant under every fixed character-corrected group action. Therefore
\[
 \|\chi(h)^{-1}h z_x\|\ge c_x
\]
for every invertible block lower triangular $h$. In particular the derivative at $h=e^{tX}$, $t=0$, is zero for each block lower triangular matrix $X$.

Put $x'=z_x/c_x$. Each normalized vector in the approximating sequence remains a product of coherent line vectors: triangular factorization of each transformed local basis expresses it as a unitary translate of the same highest line, times a scalar. The compactness of the local unitary groups gives this form also for $x'$. Moreover $Px=c_xPx'$. For a nonempty row ($r_i>0$) with vector $x'_i$, define the averaged one-slot density
\[
 A_i=\frac1{r_i}\sum_{v\text{ in row }i}
       \operatorname{Tr}_{\text{other slots}}|x'_i\rangle\langle x'_i|
     =U_i\operatorname{diag}(\theta_i)U_i^*.
\]
Here $U_i$ gives the coherent basis and $\theta_i$ its count proportions. The equality follows from the highest-count weight: the sum of diagonal occupation expectations is its count vector and all off-diagonal Lie-algebra expectations vanish. Individual slot marginals need not coincide. For an empty row choose any trace-one density; its coefficient $r_i$ is zero and it occurs in no surviving cell. The same convention applies to empty columns. The derivative just obtained says
\[
 \Re\operatorname{Tr}\!\left(X\sum_i r_iA_i\right)
 =\Re\sum_a c_{*,a}X_{aa}.
\]
Both matrices compared on the right and left are Hermitian. Testing all complex lower-triangular entries, including real diagonal entries, therefore yields
\begin{equation}\label{rec:coherent-moment-match}
 \sum_i r_iA_i=m\Lambda.
\end{equation}
Apply the same argument to $y$ to obtain $y'$, a factor $0<c_y\le1$, and $\sum_j s_jA'_j=m\Lambda$. Since $P$ is orthogonal,
\[
 |\langle y,Px\rangle|=c_xc_y|\langle y',Px'\rangle|
 \le|\langle y',Px'\rangle|.
\]
This proves both moment matching and the required coefficient compensation without assuming a finite minimizing scaling exists.

Set $Z=\Lambda^{-1/4}$, using entry one on zero coordinates, which
carry no mass.  The highest-line character gives the exact identity
\[
 Z^{\otimes m}PZ^{\otimes m}=e^{mH(\Lambda)/2}P,
 \qquad
 \log\operatorname{diag}(Z)^{c_*}=mH(\Lambda)/4.
\]
Schur orthogonality expresses $P$ as an integral of simultaneous
block-unitary actions with total absolute coefficient at most the
dimension of its compact carrier.  Insert this integral into the
identity and use the triangle inequality.  Squaring gives
\begin{equation}\label{rec:coherent-scaled-projector}
 |\langle y',Px'\rangle|^2
 \le e^{CQ^2\log(2n)-mH(\Lambda)}
       \sup_U|\langle y',(ZUZ)^{\otimes m}x'\rangle|^2.
\end{equation}
Expand each row and column vector in its coherent letter basis. Its coefficients are supported on assignments with the prescribed line counts, although they need not factor within a line. If $x_a,y_b$ are the coefficients of the resulting full assignments, then $\sum_a|x_a|^2=\sum_b|y_b|^2=1$. Cauchy--Schwarz over the pair $(a,b)$ bounds the squared matrix coefficient by the sum, over these fixed-count assignments, of the products of squared local matrix entries.

To bound that sum, sample letters independently at each cell from the row probabilities $\theta_i$ and column probabilities $\phi_j$. Every allowed full pair has the same sampling weight
$\exp\{-\sum_ir_iH(\theta_i)-\sum_js_jH(\phi_j)\}$.
Dropping the count constraints therefore bounds the last square in \eqref{rec:coherent-scaled-projector} by
\[
 e^{\sum_ir_iH(\theta_i)+\sum_js_jH(\phi_j)}
       \prod_{(i,j)\ \mathrm{survives}}W_{ij},\qquad
 W_{ij}=\Tr(A'_jZUZ A_iZU^*Z).
\]
Thus the independent sampling is a device for summing coefficients with fixed counts, not an assertion that the coherent states have identical or independent slot marginals. The matrices $A_i,A'_j$ in this formula are the averaged count densities just defined.
Under the product probabilities $r_i/m,s_j/m$, the mean of $W_{ij}$
is, by \eqref{rec:coherent-moment-match},
$\Tr(\sqrt\Lambda U\sqrt\Lambda U^*)\le1$.
The relative entropy of the uniform law on surviving cells from
the product of its marginals is at most $\log(n/m)$: compare first
with the uniform full grid and use the minimizing property of the
actual marginals.  The entropy form of the geometric-mean inequality
therefore gives
$\prod W_{ij}\le e^{m\log(n/m)}\le e^l$.
This proves \eqref{rec:coherent-coefficient}.
Use \eqref{rec:coherent-count-dimension} to pay the local entropy
factors, sum the coherent-line representations, and cancel the
global factor using $D_\mu\le e^{mH(\Lambda)}$.
Equation \eqref{rec:coherent-hook-grid} follows.
\end{proof}

By positive spectral decomposition and group Plancherel, the same
inequality gives, for group-algebra coefficients that factor over lines, $w=\bigotimes_i w_i\in\mathbb C K$ and $z=\bigotimes_j z_j\in\mathbb C J$,
\begin{equation}\label{rec:coherent-coefficient-grid}
 D_\mu\|(zw)_\mu\|_{\HS}^2
 \le e^{l+C(1+M+R)Q^2\log(2n)}
 \left(|K|\sum_{g\in K}|w_g|^2\right)
 \left(|J|\sum_{g\in J}|z_g|^2\right).
\end{equation}
The decomposition is applied to $ww^*$ and $z^*z$, so every summand
is positive.

\subsection{General diagrams by ordered holes}

We now return to the full $M$ by $R$ grid and reset the groups to
$K=S_R^M$ on rows and $J=S_M^R$ on columns.  Keeping the holes
ordered makes the intermediate list unique when the initial and
final lists are prescribed.  The resulting coefficient-square sum
has the following explicit cost.

\begin{proposition}[Restoration with ordered holes]
\label{prop:coherent-restoration}
Let $E,F$ select product unit carrier vectors of types $\sigma,\tau$
in $K,J$, respectively.  For $\lambda\vdash n$ and integer $Q\ge1$,
let $l$ be the number of boxes outside the union of its first $Q$
rows and first $Q$ columns.  Then
\begin{equation}\label{rec:coherent-general-grid}
 D_\lambda\Tr_{V_\lambda}(EF)
 \le D_\sigma D_\tau
 e^{3l+\log\binom nl+C(1+M+R)Q^2\log(2n)}.
\end{equation}
\end{proposition}

\begin{proof}
Let $\mu$ be the retained hook and $\xi$ the removed straight core.
Littlewood--Richardson inclusion places $V_\lambda$ in induction
from $V_\mu\otimes V_\xi$.
In the regular representation designate $l$ distinct labels as holes
and let $\Pi$ be the commuting relabeling projection onto type $\mu$
on the other labels.  It retains at least
$D_\mu D_\xi\ge D_\lambda/\binom nl$ dimensions in the multiplicity
space of $V_\lambda$.  Hence
\begin{equation}\label{rec:coherent-hole-lower}
 \|FE\Pi\|_{\HS}^2\ge
       \binom nl^{-1}D_\lambda\Tr_{V_\lambda}(EF).
\end{equation}

To estimate this regular-representation norm, let $\mathcal H_a$
be the span of permutations whose designated hole labels occupy
the ordered list $a$.  Fixing $a$ leaves all bijections of the
surviving labels with the $m=n-l$ other positions.  Thus
$\mathcal H_a$ is a regular $S_m$ module, with its counting basis;
$\Pi$ is the right isotypic projection onto $\mu$ on every fiber.
The row and column operators act on positions on the left.

Let $E_{ba}:\mathcal H_a\to\mathcal H_b$ and
$F_{cb}:\mathcal H_b\to\mathcal H_c$ be their blocks.  For given
$a,c$, at most one middle list $b$ is possible: each labeled
hole keeps its row under $E$ and its column under $F$.
Write $e_g,f_g$ for the group-algebra coefficients of $E,F$.
For a feasible triple choose $u\in K$, $v\in J$ with $ua=b$,
$vb=c$, and put
\[
\begin{aligned}
 K_b&=\{h\in K:hb=b\},&
 w_{ba}(h)&=e_{hu}\quad(h\in K_b),\\
 J_b&=\{t\in J:tb=b\},&
 z_{cb}(t)&=f_{vt}\quad(t\in J_b).
\end{aligned}
\]
These stabilizers fix the ordered list pointwise.  The extracted
coefficients still factor over lines.  If $L_u,L_v$ denote the
baseline motions and $W_{ba},Z_{cb}$ the remaining actions on
$\mathcal H_b$, then
$E_{ba}=W_{ba}L_u$ and $F_{cb}=L_vZ_{cb}$.
The baseline motions are unitary and commute with the right
projection $\Pi$.  Regular Plancherel on the surviving labels
therefore gives the precise block identity
\begin{equation}\label{rec:coherent-fiber-identity}
 \|F_{cb}E_{ba}\Pi|_{\mathcal H_a}\|_{\HS}^2
 =\|Z_{cb}W_{ba}\Pi|_{\mathcal H_b}\|_{\HS}^2
 =D_\mu\|(z_{cb}w_{ba})_\mu\|_{\HS}^2.
\end{equation}
The factor $D_\mu$ is the right multiplicity in the regular
$S_m$ representation.

Define the two nonnegative coefficient factors
\[
 \alpha_{ba}=|K_b|\sum_{h\in K_b}|w_{ba}(h)|^2,
 \qquad
 \beta_{cb}=|J_b|\sum_{t\in J_b}|z_{cb}(t)|^2.
\]
Equation~\eqref{rec:coherent-coefficient-grid} bounds the block
square by $e^{l+C(1+M+R)Q^2\log(2n)}\alpha_{ba}\beta_{cb}$.
For fixed $b$, the corresponding cosets partition $K$ and $J$,
so Fourier coefficient orthogonality gives
\[
 \sum_a\alpha_{ba}=\frac{|K_b|}{|K|}D_\sigma,
 \qquad
 \sum_c\beta_{cb}=\frac{|J_b|}{|J|}D_\tau.
\]
Here $\sum_{g\in K}|e_g|^2=D_\sigma/|K|$ because the selected
carrier projection has rank one, and similarly for $F$.
Since the middle list is unique, summing the block squares first
over $a,c$ and then over $b$ bounds $\|FE\Pi\|_{\HS}^2$ by
\[
 e^{l+C(1+M+R)Q^2\log(2n)}D_\sigma D_\tau
       \sum_b\frac{|K_b|}{|K|}\frac{|J_b|}{|J|}.
\]
Falling factorials in the full lines give
\[
 |K_b|/|K|\le(e/R)^l,\qquad
 |J_b|/|J|\le(e/M)^l.
\]
There are at most $n^l$ middle hole lists, and $MR=n$.
Their factors cancel to $e^{2l}$.  Add the $e^l$ in
\eqref{rec:coherent-coefficient-grid} and the binomial factor in
\eqref{rec:coherent-hole-lower} to obtain
\eqref{rec:coherent-general-grid}.
\end{proof}

\subsection{Closing the uniform moment}

\begin{proof}[Proof of Theorem~\ref{rec:coherent-main}]
Split a sweep into the balanced grid, $B_n=YX$, and suppose the
children satisfy \eqref{rec:coherent-unit-moment} with exponent $p$.
For positive $A,D$ and $p\ge1$, the Araki--Lieb--Thirring trace inequality \cite{araki1990} (see \cite[Theorem~1]{audenaert2007})
$\Tr(A^{1/2}DA^{1/2})^p\le\Tr(A^pD^p)$ gives
\begin{equation}\label{rec:coherent-trace-product}
 D_\lambda\Tr(B_n^*B_n)^p
 \le D_\lambda\Tr\{(Y^*Y)^p(XX^*)^p\}.
\end{equation}
Resolve the right side into positive sums of product projections.
Take $Q=\lfloor n^{1/50}\rfloor$ in
\eqref{rec:coherent-general-grid}.  For each row or column type,
dimension times the sum of its spectral coefficients is at most
one by the child induction.  The logarithm of the number of types
is $O(n^{3/4}\log(2n))$.  Thus
\begin{equation}\label{rec:coherent-moment-error}
 D_\lambda\Tr(B_n^*B_n)^p
 \le\exp C\{n^{4/5}+l+l\log(n/l)\}.
\end{equation}
The $Q^2$ error is also absorbed by $n^{4/5}$.  We use the
continuous convention $l\log(n/l)=0$ when $l=0$.

Put $k=n-\lambda_1$ and $L_\lambda=\log D_\lambda$. We consider nontrivial surviving types, so $k\ge1$; the trivial type will be handled separately.  The hook formula on the lower
diagram gives
\begin{equation}\label{rec:coherent-core-dimension}
 l\log Q\le L_\lambda+Ck.
\end{equation}
Indeed its dimension is at most $D_\lambda$ by branching, all its
hooks are at most $k$, and the $l$ deep-core hooks are at most $2k/Q$.
Lemma~\ref{lem:surviving-dimension} gives $L_\lambda\ge ck$
for every nontrivial surviving type.
Consequently $l\le CL_\lambda/d$.
In the nonsparse range $\log(n/k)<d^{3/4}$, this gives
\begin{equation}\label{rec:coherent-small-error}
 l\log(n/l)
 =l\log(n/k)+l\log(k/l)\le Cd^{-1/4}L_\lambda.
\end{equation}
For the second term, $k\le CL_\lambda$ and
$l\le CL_\lambda/d$ give
$l\log(k/l)\le O(L_\lambda\log d/d)$: apply monotonicity of
$x\log(CL_\lambda/x)$ on this range, increasing the constant
inside the logarithm if needed.  Also $k>ne^{-d^{3/4}}$ implies
$n^{4/5}\le Cd^{-1/4}L_\lambda$.
Consequently \eqref{rec:coherent-moment-error} is at most
$e^{C'd^{-1/4}L_\lambda}$, while it gives the operator bound
$\|B_n\|^{2p}\le e^{-(1-C'd^{-1/4})L_\lambda}$.
To see how the exponent pays the error, write
$\varepsilon=C'd^{-1/4}\le1/2$ and $T=B_n^*B_n$ on this type.
For $t\ge0$ the two estimates give
\[
 D_\lambda\Tr T^{p(1+t)}
 \le \bigl(D_\lambda\Tr T^p\bigr)\|T\|^{pt}
 \le D_\lambda^{\varepsilon-t(1-\varepsilon)}.
\]
Taking $t=2\varepsilon$ makes the exponent nonpositive.  Thus an
increase by $1+C''d^{-1/4}$ closes the estimate.

Choose an absolute base $d_0$ before choosing a sufficiently large
finite $p_0$.  On the base range, products of matching projections
have no nonconstant common fixed vector, so such $p_0$ exists.
Choose it also large enough for the sparse consequence of
\eqref{rec:coherent-contact-bound} and
\eqref{rec:coherent-tail-bound}.
Above the base take
\[
 p_d=(1+C''d^{-1/4})
       \max\{p_{\lfloor d/2\rfloor},p_{\lceil d/2\rceil}\}.
\]
Traces of positive contractions decrease with the exponent.
Thus the child estimate applies at the larger child exponent, and
the preceding argument proves the induction.  The logarithmic
increments are summable along halving scales, so $p_d\le p_*<\infty$.
The trivial representation has trace exactly one; killed types have
trace zero.  This proves \eqref{rec:coherent-unit-moment} in all cases.

The Plancherel argument of Section~\ref{ten:weighted:part:285}
applies with the dimension power $1/(2p_*)$ supplied by this
moment.  An absolute number of independent forward sweeps therefore
mixes the full permutation in $O(d)$ physical shuffles, uniformly
over the starting deck.  Proposition~\ref{ten:support} supplies
the matching lower order.
\end{proof}
